\section{A single unknown launch law from directional short flights}
\label{sec:single-law-rigidity}
We now keep one stationary launch law and allow the direction and
length of a command to be prescribed.  The output of each attempted
preparation is still the bit $B_a$ of Section~\ref{sec:setting}; neither
the actual start nor a collision location or time is observed.  The
launch density and its support are unknown.  No relation between
different launch laws is assumed, because only one law is used.

Let $\mathcal O=\bigcup_{C\in\mathcal C}C$ be a nonempty locally
finite union of compact convex bodies with nonempty interiors,
diameters at most $D$, and mutual distances at least $d>0$.
In this section each boundary is $C^2$ with positive curvature, and
every width of every body is at least $w_*>0$.  Write
\[
 n_\phi=(\cos\phi,\sin\phi),\qquad
 \tau_\phi=(-\sin\phi,\cos\phi),
\]
and denote the point with outward normal $n_\phi$ and its radius of
curvature by $c_C(\phi)$ and $r_C(\phi)$, respectively.  Thus, for the
laboratory support function $h_C$,
\begin{equation}\label{eq:single-gauss-parametrization}
 c_C(\phi)=h_C(\phi)n_\phi+h_C'(\phi)\tau_\phi,
 \qquad r_C(\phi)=h_C(\phi)+h_C''(\phi)>0.
\end{equation}
The footprint $A$ is a compact convex body with nonempty interior,
and $j$ is an $L^1$ probability density supported on $A$ and positive
almost everywhere in its interior.  Their common position is
unrestricted.  The geometric separation assumption is
\begin{equation}\label{eq:single-footprint-separation}
 \diam A\le\Delta< b_*:=\min\{d,w_*\}.
\end{equation}
For example, an inball of radius $r_*$ in each obstacle permits
$w_*=2r_*$.  No periodicity, finite species list, footprint smoothness,
density modulus of continuity, or homothety assumption is imposed.

For $0<t<t_*$ and $\theta\in\R/2\pi\Z$, the observed mean is
\begin{equation}\label{eq:single-raw-response}
 F_{t,\theta}(x)=\int_A j(z)B_{t n_\theta}(x+z)\dd z.
\end{equation}
Here $t_*>0$ is any fixed upper length; limits below use $t<d$.
All solid starts and free misses contribute zero and remain in the
denominator.  The exact datum consists of these raw forward means at
all nominal centers and all prescribed directions for arbitrarily
short positive lengths.  This direction-resolved command family is
richer than a fixed pooled compass field.  It requires no second
setting and no prescribed reciprocal joint preparation.

\begin{lemma}[The short-flight boundary measure]
\label{lem:single-boundary-flux}
As $t\downarrow0$, the fields $t^{-1}F_{t,\theta}$ converge in
$L^1_{\mathrm{loc}}(\R^2)$, locally uniformly in $\theta$, to
\begin{align}
 q_\theta(x)
 &=\sum_{C\in\mathcal C}\int_{\partial C}
       j(y-x)(-n_C(y)\cdot n_\theta)_+\dd s(y)
       \label{eq:single-flux-boundary}\\
 &=\sum_{C\in\mathcal C}\int_0^{2\pi}
       r_C(\phi)j(c_C(\phi)-x)
                 (-\cos(\theta-\phi))_+\dd\phi.
       \label{eq:single-flux-gauss}
\end{align}
For a compactly supported $C^1$ test function $\psi$, there is the
quantitative weak estimate
\begin{equation}\label{eq:single-weak-flight-error}
 \left|\int\psi(x)\{t^{-1}F_{t,\theta}(x)-q_\theta(x)\}\dd x\right|
 \le\frac t2\|\nabla\psi\|_\infty
       \sum_{C\in\mathcal C_\psi}\mathcal W_\theta(C).
\end{equation}
Here $0<t\le t_0<\min\{d,t_*\}$,
$\mathcal W_\theta(C)$ is the width perpendicular to $n_\theta$,
and $\mathcal C_\psi$ consists of the finitely many bodies meeting
$\operatorname{supp}\psi+A+t_0\overline B$.
\end{lemma}
\begin{proof}
Use coordinates $y=w\tau_\theta+u n_\theta$.  For $w$ in the
projection interval $I_{C,\theta}$, the section of $C$ is an interval
$[\ell_{C,\theta}(w),u_{C,\theta}(w)]$.  Its incoming boundary
point is
$y_{C,\theta}(w)=w\tau_\theta+\ell_{C,\theta}(w)n_\theta$.
The free starts whose segment first meets this body form precisely
the strip
\[
 \{y_{C,\theta}(w)-s n_\theta:
               w\in I_{C,\theta},\ 0<s\le t\},
\]
up to planar null sets.  Since $t<d$, a segment cannot meet two
different bodies and this strip meets no other solid.  The strips
therefore give, almost everywhere in $x$,
\begin{equation}\label{eq:single-strip-disintegration}
 \frac{F_{t,\theta}(x)}t
 =\sum_C\int_{I_{C,\theta}}\frac1t\int_0^t
   j(y_{C,\theta}(w)-s n_\theta-x)\dd s\dd w.
\end{equation}
On a compact set of nominal centers only finitely many bodies occur,
uniformly in $\theta$ and $0<t\le t_0$.  Translation continuity in
$L^1(\R^2)$ shows that replacing $s$ by zero in
\eqref{eq:single-strip-disintegration} has local $L^1$ error at most
$ND\omega_j(t)$, where $N$ is this finite component count and
\[
 \omega_j(t)=\sup_{|v|\le t}\|j(\,\cdot+v)-j\|_{L^1(\R^2)}
                         \longrightarrow0.
\]
The projection identity
$\dd w=(-n_C\cdot n_\theta)_+\dd s$ on the incoming boundary
proves \eqref{eq:single-flux-boundary}.  The Gauss parametrization
and $\dd s=r_C(\phi)\dd\phi$ give
\eqref{eq:single-flux-gauss}.

To prove the weak estimate, integrate
\eqref{eq:single-strip-disintegration} against $\psi$ and change
variables from $x$ to $z=y_{C,\theta}(w)-s n_\theta-x$.
The two test values differ by at most
$s\|\nabla\psi\|_\infty$.  Integration against the probability
density $j$ preserves this bound, and averaging $s$ over $[0,t]$
gives $t/2$.  Finally,
$|I_{C,\theta}|=\mathcal W_\theta(C)$, and bodies outside
$\mathcal C_\psi$ make no contribution.
\end{proof}

\begin{lemma}[Angular resolution of the launch density]
\label{lem:single-angular-resolution}
The field $q$ is twice continuously differentiable in $\theta$ with
values in $L^1_{\mathrm{loc}}(\R^2)$.  Its angular derivative
\begin{equation}\label{eq:single-angular-field}
 S_\theta:=(\partial_\theta^2+1)q_\theta
   =\sum_{C\in\mathcal C}\sum_{\sigma\in\{-1,1\}}
       r_C(\theta+\sigma\pi/2)
       j(c_C(\theta+\sigma\pi/2)-\,\cdot\,)
\end{equation}
is nonnegative.  For every $\theta$, the connected components of
the support of the measure $S_\theta(x)\dd x$ are exactly
\begin{equation}\label{eq:single-angular-components}
 K_{C,\sigma,\theta}=c_C(\theta+\sigma\pi/2)-A.
\end{equation}
Different such components have distance at least $b_*-\Delta>0$.
\end{lemma}
\begin{proof}
For $k(\alpha)=(-\cos\alpha)_+$, direct differentiation on the
two open semicircles and the jumps of $k'$ give, as periodic
distributions,
\begin{equation}\label{eq:single-cosine-distribution}
             k''+k=\delta_{\pi/2}+\delta_{-\pi/2}.
\end{equation}
Apply this identity to \eqref{eq:single-flux-gauss}.  Local finiteness
justifies termwise distributional differentiation.  The right side
of \eqref{eq:single-angular-field} is continuous in $\theta$ in
$L^1$ on each compact set: each $r_C$ and $c_C$ is continuous,
translations act continuously on $L^1$, and only finitely many bodies
can contribute there.  The same is true of $q$ and its first
derivative, obtained by using the bounded piecewise continuous kernel
$k'$.  The distributional identity then proves the asserted $C^2$
regularity and specifies $S_\theta$ canonically for every angle.

The support of the positive measure with density
$j(c-x)$ is exactly $c-A$.  Positivity almost everywhere in
$\operatorname{int}A$ proves this assertion even if a chosen
representative of $j$ vanishes on a dense null set.  The weights in
\eqref{eq:single-angular-field} are positive.  Contact points on
distinct obstacles have distance at least $d$.  The two contact
points belonging to one obstacle have opposite normals, and
\[
 |c_C(\phi)-c_C(\phi+\pi)|
 \ge (c_C(\phi)-c_C(\phi+\pi))\cdot n_\phi
 =h_C(\phi)+h_C(\phi+\pi)\ge w_*.
\]
For two contact points $c,c'$ and any $a,a'\in A$,
$|(c-a)-(c'-a')|\ge|c-c'|-\Delta$.
Consequently the copies in \eqref{eq:single-angular-components}
have the stated positive separation.  They are connected compact
sets and form a locally finite family.  They are therefore precisely
the connected components of the measure support.
\end{proof}

Let $s(H)$ denote the Steiner point of a convex body $H$.  Its
translation covariance and $s(-H)=-s(H)$ will fix the common
translation of the unknown table and launch law.  Define
\begin{equation}\label{eq:single-canonical-parameters}
 A_0=A-s(A),\qquad j_0(z)=j(z+s(A)),\qquad
                      \mathcal O_0=\mathcal O-s(A).
\end{equation}

\begin{theorem}[Single-law directional-germ rigidity]
\label{thm:single-law-rigidity}
Under \eqref{eq:single-footprint-separation}, the directional
short-flight germ $\{q_\theta:\theta\in\R/2\pi\Z\}$ determines
the complete canonical triple $(\mathcal O_0,A_0,j_0)$, with equality
of densities understood almost everywhere.  In particular, the raw
mean family \eqref{eq:single-raw-response} determines this triple.

More explicitly, select any angle and any connected component $K$
of $\operatorname{supp}(S_\theta(x)\dd x)$.  Let
$S_{\theta,K}=\one_KS_\theta$ and
$m_K=\int_KS_\theta(x)\dd x$.  Then
\begin{equation}\label{eq:single-density-recovery}
 A_0=-(K-s(K)),\qquad
 j_0(z)=\frac{S_{\theta,K}(s(K)-z)}{m_K}.
\end{equation}
To recover the table without matching components across angles, put
\begin{equation}\label{eq:single-boundary-recovery}
 \Gamma=\bigcup_{\theta\in\R/2\pi\Z}
       \{s(K):K\text{ is a connected component of }
                        \operatorname{supp}(S_\theta(x)\dd x)\}.
\end{equation}
Then $\Gamma=\partial\mathcal O_0$; its connected components are
the individual obstacle boundaries, and
\begin{equation}\label{eq:single-table-recovery}
 \mathcal O_0=\bigcup_{\Lambda\in\operatorname{Comp}(\Gamma)}
                                      \operatorname{conv}\Lambda.
\end{equation}
\end{theorem}
\begin{proof}
By Lemma~\ref{lem:single-angular-resolution}, the selected component
has the form $K=c-A$, and on that component the field equals
$rj(c-x)$ for one strictly positive curvature radius $r$.
Consequently
\[
 s(K)=c-s(A),\qquad m_K=r\int_Aj=r.
\]
Reflection about $s(K)$ and division by this mass give
\eqref{eq:single-density-recovery}.  This proves that the result is
independent of the selected component and angle and recovers both
the support and the complete centered density.

Again by \eqref{eq:single-angular-components}, the set of component
Steiner points at angle $\theta$ is
\[
 \{c_C(\theta-\pi/2)-s(A),
               c_C(\theta+\pi/2)-s(A):C\in\mathcal C\}.
\]
As $\theta$ ranges over the circle, the Gauss parametrization ranges
over every boundary point of every obstacle.  Thus their union is
exactly $\bigcup_C\partial(C-s(A))$.
Each of these boundaries is connected, different boundaries have
distance at least $d$, and the family is locally finite.
Their union is closed and its connected components are precisely
the individual boundaries.  The same separation shows that this
union equals the boundary of the obstacle union.  A compact convex
body with nonempty interior is the convex hull of its boundary,
which proves \eqref{eq:single-table-recovery}.
\end{proof}

\begin{corollary}[The complete observational fiber]
\label{cor:single-law-fiber}
Two configurations and single stationary launch laws in the stated
class have the same directional germ if and only if there is
$h\in\R^2$ such that
\begin{equation}\label{eq:single-law-fiber}
 \mathcal O'=\mathcal O+h,\qquad A'=A+h,
                         \qquad j'(z)=j(z-h)\quad\text{a.e.}
\end{equation}
These conditions are also equivalent to equality of the entire
raw forward response for every nominal center and every finite
command displacement for which the bit is defined.
\end{corollary}
\begin{proof}
Equal germs give equal canonical triples by
Theorem~\ref{thm:single-law-rigidity}.  Taking
$h=s(A')-s(A)$ yields \eqref{eq:single-law-fiber}.
Conversely, under these relations couple a start $x+Z$ for the first
table with $x+Z+h$ for the second.  The two unreflected segments and
all obstacles differ by the same translation.  The bits therefore
agree for every displacement, including the solid-start convention.
Equality of the full raw family implies equality of its short-flight
germ by Lemma~\ref{lem:single-boundary-flux}.
\end{proof}

For the germ, translation equality is understood in
$L^1_{\mathrm{loc}}(\R^2)$, equivalently almost everywhere in the
nominal center for each angle.  For the continuous raw mean fields
it is pointwise.  The germ period group is therefore independent
of representatives of the density and flux fields.

\begin{corollary}[Germ periods and full response completion]
\label{cor:single-germ-periods}
The first-order directional germ determines the full forward
collision law at all finite command lengths.  Moreover,
\begin{equation}\label{eq:single-germ-periods}
 \operatorname{Per}(\{q_\theta\}_\theta)
 =\operatorname{Per}(\{F_{t,\theta}\}_{0<t<t_*,\theta})
 =\operatorname{Per}(\mathcal O).
\end{equation}
The common group is discrete.  It has rank two exactly when the
configuration admits a rank-two lattice of periods with finitely
many component orbits.  No periodicity hypothesis enters this
criterion.
\end{corollary}
\begin{proof}
The canonical triple supplies the defining integral of every
forward bit, and a common translation changes none of these means.
This proves response completion.  A period of the raw family is a
period of its germ.  A period of the germ preserves each measure
support in Lemma~\ref{lem:single-angular-resolution}, permutes its
connected components, and translates their Steiner points.
It therefore preserves $\Gamma$ and, by
\eqref{eq:single-table-recovery}, the obstacle union $\mathcal O_0$.
Every period of $\mathcal O$ is conversely a period of all raw
means by the collision-bit definition.
A nonzero translation period moves any compact obstacle to a
different one, and hence has length at least $d$.  The group is
discrete; the rank-two assertion follows from local finiteness in
a compact fundamental cell, as in
Theorem~\ref{thm:global-response-rigidity}.
\end{proof}

All these conclusions concern the exact direction-resolved active
germ on the plane.  The angular operation in
\eqref{eq:single-angular-field} is applied to response means and
is not an additional sensor output.  Its estimation from finitely
many attempted bits requires regularization and separate resource
bounds.  The proof above uses neither a reciprocal occupation inverse
nor a homothetic second footprint: angular differentiation separates
the two tangent contacts, and their common density identifies the
unknown launch law.  The small-footprint condition ensures that these
translated density copies can be separated spatially; no conclusion
without that condition is asserted here.
